\section{Implementation and Result}
\label{sec:implementation}

\subsection{Implementation}

The detector of Section~\ref{sec:methodology} is deployed as the three-tier
system of Fig.~\ref{fig:f1}; the evaluation below reproduces the deployed
model service's protocol exactly, under the pinned environment (scikit-learn
1.7.2, numpy 2.2.6, pandas 2.3.3, seed $20260804$), and every value here was
re-executed by a single reproduction script against the committed data. The
primary corpus is $30{,}000$ generated rows, split $60/20/20$ in source-row
order (test positive rate $0.5027$); this is not a temporal split, since the
generator emits no timestamp, and the row-shuffled ablation of
Table~\ref{tab:t6} establishes only that row order introduces no artifact, a
weaker property than a genuine temporal holdout. Features use the shipped
scaler constants, whose legacy provenance defect is disclosed in
Section~\ref{sec:methodology} and whose cost is quantified below.

\subsection{Performance Evaluation}

The primary detector attains test AUC $0.6489$ (validation $0.6585$, indicating
no overfitting) against a chance baseline of $0.5000$ and a majority-class
accuracy of $0.5027$; all other measures (PR-AUC, precision, recall, $F_1$,
MCC, balanced accuracy) are in Table~\ref{tab:t6} and each exceeds its
baseline. The probability output is usable as a probability, at Brier score
$0.2336$ (baseline $0.25$) and expected calibration error $0.0214$ over ten
bins. Five-fold
cross-validation gives AUC $0.6519\pm0.0094$ and PR-AUC $0.6452\pm0.0065$, and
a bootstrap over $1{,}000$ test resamples gives a $95\%$ CI of
$[0.6342,0.6622]$ for test AUC, excluding both the majority-class rate and the
random forest's $0.6790$ below, so the margin over chance is not an artifact
of the test split. At the deployed threshold $0.5$, recall is a conservative
$0.5212$; sweeping the threshold trades this against precision as expected
(recall $0.9907$ at $t{=}0.3$, precision $0.7675$ at $t{=}0.7$), so a
deployment with asymmetric costs should choose its operating point rather than
inherit $0.5$.

Table~\ref{tab:t6} compares the primary model with alternatives on the
identical split: regularization strength is inert ($C\in\{0.1,1,10\}$ all give
test AUC $0.6489$), standardization is not (raw features fall to $0.6333$),
and the best measured model is a random forest at $0.6790$, followed by
gradient boosting ($0.6743$) and a calibrated RBF SVM ($0.6702$) --- a
deficit of $0.0301$ AUC that the logistic detector accepts for the
aggregability and calibration argued in Section~\ref{sec:methodology}.

\subsection{Findings and Comparison with Existing Approaches}

The central finding is a matched negative control that varies the label
function and nothing else. The same $30{,}000$ feature rows are relabelled by
the legacy heuristic scoring rule in place of the backtested
certainty-equivalent criterion~\eqref{eq:e23}; rows, features, standardization,
split, model class and evaluation code are identical across both arms.
Fig.~\ref{fig:f11} shows the result: test AUC $0.9890$ under heuristic labels
against $0.6489$ under backtested labels, a gap of $0.3401$ attributable to
the label alone. The mechanism is directly observable: elapsed time correlates
$0.6074$ with the heuristic label, which reads it as an input, and $0.0256$
with the backtested label, which contains no elapsed-time term; the two
labellings agree on only $0.5221$ of rows. The figure $0.9890$ is a degeneracy
exhibit, not an achievement --- it is the score returned when the oracle
scored against is the rule the model was given to imitate --- and $0.6489$ is
the honest difficulty of the task once that oracle is replaced by a forward,
net-of-cost outcome. A separate $5{,}000$-row corpus labelled by the same
heuristic reaches test AUC $0.9310$ and is reported only as corroboration,
since its corpus and class balance differ and the shipped scaler is in-sample
for it.

Table~\ref{tab:t7} scores the conventional policies against outcome labels on
the same test rows. The calendar rule attains AUC $0.5167$, the band trigger
$0.5146$, and --- the argumentative payload --- the legacy heuristic applied
directly as a decision rule attains only $0.5278$: the rule that a naively
supervised model reproduces at $0.9890$ predicts actual net-of-cost benefit
barely above chance, so recovering the rule and deciding well are almost
unrelated objectives. Paired McNemar tests confirm the ordering in both
directions (Table~\ref{tab:t7}): the detector significantly beats every
deterministic rule and the heuristic itself ($p<10^{-16}$ throughout), and is
itself significantly beaten by both tree ensembles. Against the alternative
learners the comparison runs the other way and is reported as measured: the
random forest's $0.0301$ AUC advantage is not within noise, but its
parameters admit no meaningful coordinate-wise average, whereas the linear
model's twelve-weight-plus-intercept vector is precisely what federated
averaging operates on. The trade is worth making only because the federated
path is pursued; were federation abandoned, a non-federated nonlinear model
would be the correct choice on this data.

The fitted coefficients (Table~\ref{tab:t2}) are dominated by
\texttt{market\_volatility\_30d} ($+1.0835$) and \texttt{portfolio\_return}
($+1.0661$), following directly from the structure of~\eqref{eq:e23}, while
\texttt{top3\_concentration} carries a negative coefficient despite a positive
univariate correlation --- the signature of a concentration family correlated
at $0.91$--$0.92$. No economic reading is claimed for the small coefficients.

One caveat applies throughout: $0.6489$ is the honest difficulty of the task
under outcome labels, not a ceiling --- no Bayes-rate estimate was computed, and
Table~\ref{tab:t6} already shows nonlinear learners exceeding it, so the claim
is task hardness, not optimality. This evaluation reports classification
metrics only: no path-based backtest was performed, so no Sharpe ratio,
Sortino ratio, maximum drawdown, turnover or annualised return is reported
anywhere in this paper, and no federated experiment was run --- both are
future work.
